% Part: second-order-logic
% Chapter: sol-and-sets
% Section: introduction

\documentclass[../../../include/open-logic-section]{subfiles}

\begin{document}

\olfileid{sol}{set}{int}

\olsection{سريزه}

څرنګه چې دويمې درجې منطق د تعريف ساحې پر فرعي سټونو او هم پر
تابعو کميت اچولے شي، تمه کېږي چې د سټونو د تيورۍ لږ تر لږه يوه
برخه په دې منطق کښې بيانېدے شي۔ له «بيانېدو» څخه مو مطلب دا دے
چې د سټونو د تيورۍ خاصيتونه او ويناوې په دويمې درجې منطق کښې،
بې له دې چې د سټونو دپاره کوم ځانګړے غيرمنطقي لغت وکاروو
(لکه د سټونو د تيورۍ د غړيتوب !!{predicate})، څرګندولے شو۔
د بېلګې په توګه، د سټونو اتحاد، اشتراک او د فرعي سټ اړيکه
تعريفولے شو؛ د سټونو اندازې هم سره پرتله کولے شو او د
کانتور د قضيې غوندې پايلې بيانولے شو۔

\end{document}
